
\documentclass{article}
\usepackage{colortbl}
\usepackage{makecell}
\usepackage{multirow}
\usepackage{supertabular}

\begin{document}

\newcounter{utterance}

\centering \large Interaction Transcript for game `wordle', experiment `medium\_frequency\_words\_no\_clue\_no\_critic', episode 10 with reflexion{-}Qwen3.5{-}9B.
\vspace{24pt}

{ \footnotesize  \setcounter{utterance}{1}
\setlength{\tabcolsep}{0pt}
\begin{supertabular}{c@{$\;$}|p{.15\linewidth}@{}p{.15\linewidth}p{.15\linewidth}p{.15\linewidth}p{.15\linewidth}p{.15\linewidth}}
    \# & \multicolumn{2}{c}{Player} && \multicolumn{2}{c}{Game Master} \\
    \hline

    \theutterance \stepcounter{utterance}  
    & & & \multicolumn{4}{p{0.6\linewidth}}{
        \cellcolor[rgb]{0.9,0.9,0.9}{
            \makecell[{{p{\linewidth}}}]{
                \texttt{\tiny{[P1$\langle$GM]}}
                \texttt{You have attempted this same task before and did not succeed. These are the plans you wrote after those attempts.} \\
\texttt{Trial \#1: New plan: I failed because I guessed "tent" (4 letters) instead of a 5{-}letter word, violating the game rules. The feedback "t<green> e<green> n<green> t<green> s<red>" actually indicated that the target word starts with "tent" followed by an incorrect letter 's', meaning the word is likely "tents". However, my explanation incorrectly claimed "tent" was the full target word, leading to an invalid 4{-}letter guess. Next time, I will carefully count the letters in my guess to ensure it is exactly five characters long, and I will verify that my explanation accurately reflects the state of the game before submitting. I should have recognized that "tent" + 's' = "tents" and guessed "tents" immediately after the previous turn.} \\
\\ 
\texttt{Trial \#2: New plan: I failed because I got stuck in a loop of internal monologue, repeatedly rejecting "enhance" (which is 7 letters) and failing to generate a valid 5{-}letter word that fits the pattern `\_e\_\_h` with 'n' included. I violated the rule "Do not generate any other text" by outputting a long stream of reasoning instead of a concise explanation followed immediately by a valid guess. Next time, upon receiving feedback, I will immediately construct a valid 5{-}letter word (e.g., "neigh" was good, but I need to solve `\_e\_n\_h` or similar) that satisfies the constraints, ensure it is exactly five characters long, and output only the required "explanation:" and "guess:" tags without any extra commentary or hesitation.} \\
\\ 
\texttt{Trial \#3: New plan: I failed because I got stuck in a loop of internal monologue, repeatedly rejecting valid candidates like "enhance" (which is 7 letters) and failing to generate a valid 5{-}letter word that fits the pattern `\_e\_n\_h` with 'n' included. I violated the rule "Do not generate any other text" by outputting a long stream of reasoning instead of a concise explanation followed immediately by a valid guess. Next time, upon receiving feedback, I will immediately construct a valid 5{-}letter word (e.g., "neigh" was good, but I need to solve `\_e\_n\_h` or similar) that satisfies the constraints, ensure it is exactly five characters long, and output only the required "explanation:" and "guess:" tags without any extra commentary or hesitation.} \\
\\ 
\\ 
\texttt{You are a language wizard who likes to guess words by using the given rules.} \\
\\ 
\texttt{Welcome to Wordle! You have six attempts to guess the target word, a valid English word of five lowercase letters (a{-}z). Please use the tags "explanation:" and "guess:" to provide a concise explanation for each guess.} \\
\\ 
\texttt{For instance, if your guess is "apple", your response should be} \\
\texttt{explanation: this is a common five{-}letter English word, and I am starting my guess with this word.} \\
\texttt{guess: apple} \\
\\ 
\texttt{After each guess, your answer will be validated, and you will receive feedback indicating which letters are correct (green), which letters are correct but in the wrong position (yellow), and which letters are incorrect (red). This feedback can be useful in determining which letters to include or exclude in your next guess.} \\
\\ 
\texttt{For example, the feedback for "apple" might be:} \\
\texttt{guess\_feedback: a<yellow> p<yellow> p<green> l<yellow> e<red>} \\
\\ 
\texttt{The explanation should contain details about how the guess\_feedback is used to arrive at a new guess.} \\
\\ 
\texttt{Do not generate any other text.} \\
\\ 
\texttt{Let's begin with your first guess.} \\
            }
        }
    }
     \\ \\

    \theutterance \stepcounter{utterance}  
    & \multicolumn{4}{p{0.6\linewidth}}{
        \cellcolor[rgb]{0.561,0.737,0.561}{
            \makecell[{{p{\linewidth}}}]{
                \texttt{\tiny{[P1$\rangle$GM]}}
                \texttt{explanation: Starting with a common five{-}letter word to gather initial information about letter placement and validity.} \\
\texttt{guess: apple} \\
            }
        }
    }
    & & \\ \\

    \theutterance \stepcounter{utterance}  
    & & & \multicolumn{2}{p{0.3\linewidth}}{
        \cellcolor[rgb]{0.9,0.9,0.9}{
            \makecell[{{p{\linewidth}}}]{
                \texttt{\tiny{[GM$|$GM]}}
                \texttt{attempts = 1} \\
\texttt{target\_word = tenth} \\
\texttt{guess = apple} \\
\texttt{guess\_feedback = a<red> p<red> p<red> l<red> e<yellow>} \\
            }
        }
    }
    & & \\ \\

    \theutterance \stepcounter{utterance}  
    & & & \multicolumn{4}{p{0.6\linewidth}}{
        \cellcolor[rgb]{0.9,0.9,0.9}{
            \makecell[{{p{\linewidth}}}]{
                \texttt{\tiny{[P1$\langle$GM]}}
                \texttt{guess\_feedback: a<red> p<red> p<red> l<red> e<yellow>} \\
\\ 
\texttt{Provide your response only in this format.} \\
\texttt{explanation: details} \\
\texttt{guess: word} \\
            }
        }
    }
     \\ \\

    \theutterance \stepcounter{utterance}  
    & \multicolumn{4}{p{0.6\linewidth}}{
        \cellcolor[rgb]{0.561,0.737,0.561}{
            \makecell[{{p{\linewidth}}}]{
                \texttt{\tiny{[P1$\rangle$GM]}}
                \texttt{explanation: The feedback indicates that 'a', 'p', and 'l' are incorrect, while 'e' is correct but in the wrong position (not at the end). I will try a new word that includes 'e' in a different position and avoids the incorrect letters.} \\
\texttt{guess: there} \\
            }
        }
    }
    & & \\ \\

    \theutterance \stepcounter{utterance}  
    & & & \multicolumn{2}{p{0.3\linewidth}}{
        \cellcolor[rgb]{0.9,0.9,0.9}{
            \makecell[{{p{\linewidth}}}]{
                \texttt{\tiny{[GM$|$GM]}}
                \texttt{attempts = 2} \\
\texttt{target\_word = tenth} \\
\texttt{guess = there} \\
\texttt{guess\_feedback = t<green> h<yellow> e<yellow> r<red> e<red>} \\
            }
        }
    }
    & & \\ \\

    \theutterance \stepcounter{utterance}  
    & & & \multicolumn{4}{p{0.6\linewidth}}{
        \cellcolor[rgb]{0.9,0.9,0.9}{
            \makecell[{{p{\linewidth}}}]{
                \texttt{\tiny{[P1$\langle$GM]}}
                \texttt{guess\_feedback: t<green> h<yellow> e<yellow> r<red> e<red>} \\
\\ 
\texttt{Provide your response only in this format.} \\
\texttt{explanation: details} \\
\texttt{guess: word} \\
            }
        }
    }
     \\ \\

    \theutterance \stepcounter{utterance}  
    & \multicolumn{4}{p{0.6\linewidth}}{
        \cellcolor[rgb]{0.561,0.737,0.561}{
            \makecell[{{p{\linewidth}}}]{
                \texttt{\tiny{[P1$\rangle$GM]}}
                \texttt{explanation: The feedback shows 't' is correct and in the first position, 'h' and 'e' are correct but misplaced (so 'e' is not in the 3rd or 5th position), and 'r' is incorrect. Since 'e' is already known to be in the word but not at positions 3 or 5, and 't' is at position 1, I need a word starting with 't', containing 'h' and 'e' in other positions. A strong candidate is "the" related words or similar structures. Let's try "these" but wait, 'e' cannot be at 3 or 5. Actually, "these" has 'e' at 2 and 5. If 'e' is yellow in "there" at pos 3 and 5, it means 'e' is in the word but not at 3 or 5. So in "these", 'e' at 5 would be red if it's not allowed. Let's reconsider. Maybe "thems"? No, 5 letters. "theme"? t{-}h{-}e{-}m{-}e. 'e' at 3 and 5. That might be wrong. How about "thine"? t{-}h{-}i{-}n{-}e. 'e' at 5. If 'e' is not at 5, that's bad. Wait, the feedback for "there" was t<green> h<yellow> e<yellow> r<red> e<red>. The last 'e' is red, meaning 'e' is not at position 5. The middle 'e' is yellow, meaning 'e' is in the word but not at position 3. So 'e' must be at position 2 or 4. Let's try "thems" {-}> no 'e'. "theme" {-}> e at 3 and 5, both likely wrong positions. "thine" {-}> e at 5, likely wrong. "thigh"? t{-}h{-}i{-}g{-}h. No 'e'. "thief"? t{-}h{-}i{-}e{-}f. 'e' at 4. That fits! 't' at 1 (green), 'h' at 2 (yellow, but in "thief" it's at 2, so if it's yellow in "there" at 2, it means 'h' is correct but not at 2? Wait, in "there", 'h' is at 2 and it's yellow, meaning 'h' is in the word but not at 2. Oh! I misread. In "there": t<green> (pos 1 correct), h<yellow> (pos 2 wrong, but 'h' is in word), e<yellow> (pos 3 wrong, but 'e' is in word), r<red>, e<red> (pos} \\
            }
        }
    }
    & & \\ \\

    \theutterance \stepcounter{utterance}  
    & & & \multicolumn{2}{p{0.3\linewidth}}{
        \cellcolor[rgb]{0.9,0.9,0.9}{
            \makecell[{{p{\linewidth}}}]{
                \texttt{\tiny{[GM$|$GM]}}
                \texttt{Error: The response should contain the 'guess:' keyword exactly once.} \\
            }
        }
    }
    & & \\ \\

    \theutterance \stepcounter{utterance}  
    & & & \multicolumn{2}{p{0.3\linewidth}}{
        \cellcolor[rgb]{0.9,0.9,0.9}{
            \makecell[{{p{\linewidth}}}]{
                \texttt{\tiny{[GM$|$GM]}}
                \texttt{game\_result = ABORT} \\
            }
        }
    }
    & & \\ \\

    \theutterance \stepcounter{utterance}  
    & & \multicolumn{4}{p{0.6\linewidth}}{
        \cellcolor[rgb]{0.95,0.95,0.95}{
            \makecell[{{p{\linewidth}}}]{
                \texttt{\tiny{[P1$\leftrightarrow$P]}}
                \texttt{Post{-}episode Reflexion (not sent to GM):} \\
\texttt{New plan: I failed because I violated the strict output format by generating a long stream of internal monologue and reasoning instead of the required concise "explanation:" followed immediately by "guess:". The game rules explicitly state "Do not generate any other text," yet I outputted multiple paragraphs of thought process before attempting a guess. Next time, upon receiving feedback, I will immediately synthesize the constraints (e.g., 't' at pos 1, 'h' and 'e' misplaced, 'r' excluded) into a single, valid 5{-}letter word (like "thems" or "theme" if valid, or "thief") and output only the two required lines: a brief explanation of the logic and the guess itself, with no extra commentary or hesitation.} \\
            }
        }
    }
    & \\ \\

\end{supertabular}
}

\end{document}
